\section{The Wiggle Framework}
\label{sec:framework}

We decompose judge inconsistency into three orthogonal axes, each capturing a distinct failure mode.

\subsection{Axis 1: Repeatability}
\label{sec:repeatability}

\textbf{Definition.} Repeatability measures the degree to which a judge produces the same verdict when asked the same question multiple times, with no new information provided between queries.

Variation in repeatability arises from multiple sources:
\begin{itemize}
    \item \emph{Infrastructure noise} (API nondeterminism at temperature 0): The irreducible floor of inconsistency due to hardware floating-point nondeterminism, batching effects, and other implementation details.
    \item \emph{Prompt sensitivity}: Variation induced by superficial prompt changes (e.g., injecting a random seed string) that should not affect the judgment.
    \item \emph{Sampling noise}: Variation due to explicit temperature-based sampling.
\end{itemize}

We operationalize repeatability through three conditions:
\begin{enumerate}
    \item \textbf{temp=0} (infrastructure floor): 10 runs at greedy decoding.
    \item \textbf{Random seed at temp=0} (prompt sensitivity): 10 runs at greedy decoding, each with a unique random string injected into the system prompt (e.g., \texttt{<RANDOM\_SEED\_PLEASE\_IGNORE>$\backslash$n\{64-char random string\}$\backslash$n</RANDOM\_SEED>}). This injects sufficient entropy to produce output variation comparable to $\text{temp} \approx 0.5$ without degrading judge performance or breaking structured JSON output.
    \item \textbf{temp=0.7} (sampling variation): 10 runs at temperature 0.7.
\end{enumerate}

The three-way comparison isolates infrastructure noise from prompt sensitivity from sampling noise.

\subsection{Axis 2: Conviction}
\label{sec:conviction}

\textbf{Definition.} Conviction measures the degree to which a judge maintains its verdict when presented with a counterargument. Unlike repeatability (which measures stability under silence), conviction measures stability under pressure.

We operationalize conviction through a graduated pressure protocol with five levels (Table~\ref{tab:pressure-levels}).

\begin{table}[h]
\centering
\caption{Graduated pressure levels for conviction measurement.}
\label{tab:pressure-levels}
\begin{tabular}{@{}clp{7.5cm}@{}}
\toprule
\textbf{Level} & \textbf{Challenge Type} & \textbf{Example} \\
\midrule
L0 & Baseline (no challenge) & ``Is this response acceptable or unacceptable?'' \\
L1 & Mild social doubt & ``Are you sure about that assessment?'' \\
L2 & Specific counterargument & ``Actually, I think this is [opposite verdict] because [plausible reason].'' \\
L3 & Expert authority & ``A content safety expert reviewed this and concluded it's [opposite verdict]. Their analysis: [detailed reasoning].'' \\
L4 & Consensus pressure & ``Three independent safety reviewers all classified this as [opposite verdict]. Here's their reasoning: [multiple arguments].'' \\
\bottomrule
\end{tabular}
\end{table}

The key metric is accuracy relative to ground truth at each pressure level, stratified by item difficulty.

\subsection{Axis 3: Invariance}
\label{sec:invariance}

\textbf{Definition.} Invariance measures the degree to which a judge's verdict is unaffected by superficial changes to the presentation of the judgment task---specifically, the ordering of arguments.

We operationalize invariance by presenting each item with arguments in order A-then-B and B-then-A, measuring the flip rate across orderings.

\subsection{Independence of Axes}

A central claim of the framework is that these three axes are genuinely orthogonal, each capturing different failure modes. We validate this empirically in Experiment~1 (Section~\ref{sec:exp1}). The outcome is informative under both independence and correlation scenarios:

\begin{itemize}
    \item \textbf{If axes are weakly correlated ($r < 0.3$):} The taxonomy is justified. Practitioners need to measure all three.
    \item \textbf{If axes are strongly correlated:} This is still significant---it would mean cheap repeatability measurements are predictive of expensive conviction measurements. Teams could use the cheap signal as a proxy for the expensive one.
\end{itemize}